\chapter*{The Argument Before the Field}
\addcontentsline{toc}{chapter}{The Argument Before the Field}
\markboth{The Argument Before the Field}{The Argument Before the Field}

In 1766, a respected philosopher in Königsberg sat down to write about a man he could not stop thinking about.

The man was Emanuel Swedenborg, a former scientist and engineer who had spent his early career on canals, mining machinery, and treatises on metallurgy and cosmology. In midlife, something changed. Swedenborg began reporting visions: encounters with angels, journeys through spiritual worlds, and most famously, a claim that he had witnessed a fire in Stockholm from three hundred kilometers away, in real time, while sitting at a dinner party in Gothenburg. Word of the episode spread across Europe. It reached Immanuel Kant.

Kant was intrigued, then troubled. He bought Swedenborg's sprawling eight-volume \emph{Arcana Coelestia} at considerable expense and read it looking for evidence, one way or the other, of whether the visions were real. What he found instead was something harder to categorize: a system. Swedenborg's spiritual world was not a jumble of disconnected hallucinations. It was internally coherent, richly detailed, and endlessly extensible. Every new vision fit the architecture of the ones before it. Correspondences bred further correspondences. Nothing in it was obviously self-contradictory. And none of it could be checked against anything outside itself.

Kant's response, \emph{Dreams of a Spirit-Seer}, is a strange piece of writing. On the surface it mocks Swedenborg. Underneath, it is doing something more careful, working toward a single question: when does internally coherent reasoning count as knowledge, rather than merely fluent invention? That question did not leave Kant. Fifteen years later it resurfaced, transformed, as the \emph{Critique of Pure Reason}, the book that tried to specify, once and for all, the conditions under which a claim to knowledge is entitled to be believed.

This book is not about Kant or Swedenborg. It is about machine learning. But it opens here because the question they were fighting over, in 1766, is the same question that has driven nearly every major development in this field, right up to the present: when does internally coherent reasoning count as knowledge?

\section*{Two Ways to Go Wrong}

Any system capable of producing rich, structured output faces two very different failure modes.

The first is a system that cannot produce enough. It is too constrained, too literal, too dependent on rules someone wrote down in advance. It fails by having nothing to say about situations its designer did not anticipate. Symbolic AI's expert systems, built from thousands of hand-coded if-then rules, failed this way. They were precise and they were brittle. The world kept presenting cases the rules did not cover.

The second failure mode is the opposite, and it is the one Kant was staring at when he read the \emph{Arcana Coelestia}. It is a system that can produce almost anything: fluent, internally consistent, richly elaborated, and only loosely tethered to anything outside itself. Give it a pattern and it will extend that pattern indefinitely, whether or not the extension is true. Modern language models, when they hallucinate, are doing something structurally similar to what Kant suspected Swedenborg's imagination was doing: sampling forward from an internal model of what a plausible next thing looks like, unconstrained by an external check on whether that next thing actually happened.

Swedenborg was not lying, in the ordinary sense. He was, by every account, sincere. His visions were the output of an extraordinarily generative mind operating without a verification loop back to shared, external reality. That is the useful, non-mystical way to read him: not as a fraud and not as a genuine seer, but as an early, human-scale example of what unconstrained generation looks like when it is fluent enough to be convincing.

Kant's answer was not to demand more rules, in the manner of the symbolic systems that would fail two centuries later. His answer was to ask a different kind of question entirely: not "is this specific claim true," but "what are the general conditions any claim has to meet, before it is even eligible to count as knowledge." Space, time, and causality, in his account, are not things we find in the world. They are the structure a mind imposes on raw experience in order to make experience navigable at all. Reason, for Kant, is not a longer list of facts. It is an architecture that decides which internally coherent stories are allowed to stand.

\section*{Why This Belongs in a History of Machine Learning}

Every architecture discussed in this book, from the perceptron to the transformer, sits somewhere on the line between these two failure modes. Each one solves a specific bottleneck: a limit on what the previous generation of models could represent, learn, remember, compute, or verify. And nearly every solution to one bottleneck creates or exposes another, in the same way that Kant's boundary-drawing solved the problem of unconstrained metaphysical speculation while creating a new, harder problem: how do you build a system expressive enough to be useful, without it also being expressive enough to make anything sound plausible. Read that pattern carefully enough, across enough chapters, and a stronger claim starts to look unavoidable: every mechanism that increases what a system can express eventually creates a new problem of whether that expression can be trusted. This book's last chapter returns to that claim directly, once there are sixteen chapters of evidence behind it instead of one historical argument.

That tension, expressiveness against grounding, generative capacity against verification, is not a modern invention arriving with GPUs and internet-scale text corpora. It is the oldest problem in the study of minds, artificial or otherwise. Kant and Swedenborg give it a human face and a specific date: an argument conducted in letters and treatises across the 1760s and 70s, between a man who saw too much and a man who was afraid of seeing anything he could not later justify.

The chapters that follow trace how the field of machine learning has, generation after generation, reinvented versions of that argument: sometimes as a fight between symbolic rules and statistical learning, sometimes as a fight between raw generative capacity and the alignment techniques meant to constrain it, sometimes as a quieter, more technical question about whether a model's internal representation of the world is answerable to anything outside itself. The specific names change. The underlying question, the one Kant could not put down after reading Swedenborg, does not.
